\clearpage\thispagestyle{empty}
\pdfbookmark[0]{Quellen zur Besprechung der Connaissance des tems für 1811}{connaissance1811-sources}
\begin{center}\textbf{Quellen zur Besprechung der Connaissance des tems für 1811}\end{center}
\noindent Diese KI-übersetzte Fassung in modernem Deutsch folgt Carl Friedrich Gauß, \textit{Werke}, Band~VI (Göttingen, 1874), S.~527--528. Die Besprechung erschien zuerst ohne Verfassernamen unter Paris in den \textit{Göttingischen gelehrten Anzeigen}, Stück 153, 24. September 1810, S.~1526--1528. Das Bureau des longitudes gab das Jahrbuch für 1811 heraus.
\par\smallskip\noindent Gauß’ \hyperlink{native-44}{Programm zur Bestimmung der Polhöhe} und die hier angeführte \hyperlink{native-73}{Anzeige von 1808} stehen weiter vorn in diesem Band. Der Seitenverweis im Text folgt der Werkausgabe.
\par\smallskip\noindent \href{https://gdz.sub.uni-goettingen.de/id/PPN236007467}{Text der Werkausgabe}. Erstdruck, \href{https://images.p.sub.uni-goettingen.de/iiif/image/gdz:PPN319721507_1810:00001560/full/1650,/0/default.jpg}{S.~1526}, \href{https://images.p.sub.uni-goettingen.de/iiif/image/gdz:PPN319721507_1810:00001561/full/1650,/0/default.jpg}{S.~1527}, \href{https://images.p.sub.uni-goettingen.de/iiif/image/gdz:PPN319721507_1810:00001562/full/1650,/0/default.jpg}{S.~1528}. \href{https://cdt.imcce.fr/items/show/160}{Katalogeintrag zum Jahrbuch}.
\par\smallskip\noindent Übersetzung und Satz mit KI erstellt.
\clearpage
